% Slides in the Laboratory house style: 16:9 Beamer with the document faces.
% One idea per frame, a title that states it, and a figure drawn for the slide.
% Replace the pendulum with the talk and keep the conventions. Build with
% build.py, which reads latexmkrc.
\documentclass[aspectratio=169]{beamer}
\usetheme{Laboratory}

% Titles become PDF bookmarks, which hold no maths: give a plain-text twin.
\title{Measuring \texorpdfstring{$g$}{g} with a simple pendulum}
\author{Author Name}
\date{\today}
% The repository, notebook or data folder behind the talk; delete to omit.
\source{analysis/pendulum.py}

\begin{document}

\begin{frame}[plain]
  \titlepage
\end{frame}

% A frame title is the claim the frame supports, not a topic label.
\begin{frame}{The period of a small swing depends only on length}
  \begin{itemize}
    \item Released from rest at a small angle, a pendulum of length $L$ swings with period
      \begin{equation*}
        T = 2\pi \sqrt{L/g}
      \end{equation*}
    \item Timing $T$ at several lengths measures $g$ and tests the model~\cite{nelson1986}
    \item At the \qty{3}{\degree} release used here, the finite amplitude lengthens $T$ by \num{1.7e-4}
  \end{itemize}
\end{frame}

\begin{frame}{The small-angle model fits every length to within the timing error}
  \includegraphics[width=\linewidth]{figures/pendulum}

  Period against length with the fitted model (a), and residuals from the fit (b), at nine lengths timed to \qty{1}{\milli\second}.
\end{frame}

\begin{frame}{The fit gives \texorpdfstring{$g$}{g} to four parts in ten thousand}
  % On a slide a table sits beside its words, so it needs no float, number or reference.
  \begin{center}
    \begin{tabular}{
      S[table-format=1.1]
      S[table-format=1.4(2)]
      S[table-format=1.4]
    }
      \toprule
      {\colhead{$L$}{\metre}} & {\colhead{$T$}{\second}} & {\colhead{$T_\mathrm{fit}$}{\second}} \\
      \midrule
      0.2 & 0.8983(10) & 0.8972 \\
      0.6 & 1.5548(10) & 1.5540 \\
      1.0 & 2.0057(10) & 2.0061 \\
      \bottomrule
    \end{tabular}
  \end{center}
  \begin{block}{Result}
    $g = \qty{9.809(4)}{\metre\per\second\squared}$, with $\chi^2 = 6.9$ for 8 degrees of freedom~\cite{taylor1997}
  \end{block}
\end{frame}

\begin{frame}{References}
  \bibliographystyle{apsrev4-2}
  \bibliography{refs}
\end{frame}

\end{document}
